\documentclass[11pt]{article}
\usepackage{amsmath, amssymb, amsthm}
\usepackage[margin=1in]{geometry}
\usepackage{hyperref}

\newtheorem{theorem}{Theorem}
\newtheorem{proposition}{Proposition}
\newtheorem{corollary}{Corollary}
\newcommand{\bmat}[1]{\begin{pmatrix}#1\end{pmatrix}}
\newcommand{\conj}[1]{\overline{#1}}
\newcommand{\hc}{\dagger}
\newcommand{\Hhat}{\widehat{H}}
\newcommand{\reals}{\mathbb{R}}
\newcommand{\complexes}{\mathbb{C}}

\title{Structure-Preserving Diagonalization of the BSE Hamiltonian:\\
  Reduction to Hermitian and Skew-Symmetric Eigenvalue Problems}
\author{GPAW Technical Notes}
\date{\today}

\begin{document}
\maketitle

\begin{abstract}
  We present a self-contained derivation of the structure-preserving
  diagonalization method for the Bethe--Salpeter equation (BSE)
  Hamiltonian, as implemented in GPAW.  The pseudo-Hermitian structure
  of the BSE matrix is exploited to reduce the eigenvalue problem to a
  standard Hermitian eigenproblem (complex case) or a real
  skew-symmetric eigenproblem (real case), both of which can be solved
  more efficiently than the original non-Hermitian problem.  We prove
  the equivalence of eigenvalues, derive the eigenvector
  back-transformation, and show how left eigenvectors are obtained
  directly from right eigenvectors via the BSE structure.
\end{abstract}

\tableofcontents

%======================================================================
\section{The BSE Hamiltonian and its structure}
%======================================================================

The discretized Bethe--Salpeter equation gives rise to an eigenvalue
problem
\begin{equation}\label{eq:bse}
  H x = \omega x
\end{equation}
where the BSE Hamiltonian $H \in \complexes^{2n \times 2n}$ has the
block structure
\begin{equation}\label{eq:H}
  H = \bmat{A & B \\ -B^* & -A^*}
\end{equation}
with the following symmetry properties:
\begin{itemize}
  \item $A \in \complexes^{n\times n}$ is \textbf{Hermitian}: $A^\hc = A$,
  \item $B \in \complexes^{n\times n}$ is \textbf{complex symmetric}: $B^T = B$.
\end{itemize}
Here $A$ is the resonant block (containing single-particle transition
energies and direct/exchange interactions between resonant
electron--hole pairs), and $B$ is the coupling block between resonant
and anti-resonant transitions.

The indices range over pairs of electron--hole states $(k, v, c)$
comprising a crystal momentum $k$, a valence band $v$, and a
conduction band $c$. For a system with $n$ resonant pairs, the matrix
is $2n \times 2n$.

\medskip

\noindent\textbf{Key property:} The matrix $H$ is
\emph{pseudo-Hermitian}; there exists a Hermitian involutory matrix
$S$ such that $H^\hc = S H S^{-1}$, where
\begin{equation}\label{eq:S}
  S = \bmat{I_n & 0 \\ 0 & -I_n}, \qquad S = S^\hc = S^{-1}.
\end{equation}

This is easily verified:
\[
  SHS = \bmat{I & 0 \\ 0 & -I}\bmat{A & B \\ -B^* & -A^*}\bmat{I & 0 \\ 0 & -I}
  = \bmat{A & -B \\ B^* & -A^*} = H^\hc.
\]

%======================================================================
\section{The Hermitian companion matrix}
%======================================================================

Define the \emph{Hermitian companion} of $H$ by
\begin{equation}\label{eq:Hhat}
  \Hhat := S H = \bmat{A & B \\ B^* & A^*}.
\end{equation}
We verify that $\Hhat$ is Hermitian:
\begin{align*}
  \Hhat^\hc
  &= \bmat{A^\hc & (B^*)^\hc \\ B^\hc & (A^*)^\hc} \\
  &= \bmat{A & \conj{B^*}^T \\ \conj{B}^T & \conj{A^*}^T}.
\end{align*}
Using $A^\hc = A$ (Hermitian) and $B^T = B$ (symmetric):
\begin{itemize}
  \item $\conj{B^*}^T = B^T = B$, so the $(1,2)$ block is $B$. \checkmark
  \item $\conj{B}^T = \conj{B^T} = \conj{B} = B^*$, so the $(2,1)$
    block is $B^*$. \checkmark
  \item For the $(2,2)$ block: since $A^\hc = A$ implies $A^T = \conj{A}$,
    we have $\conj{A^*}^T = A^T = \conj{A} = A^*$. \checkmark
\end{itemize}
Hence $\Hhat^\hc = \Hhat$.

\begin{equation}
  \boxed{H = S \Hhat, \qquad \Hhat = \Hhat^\hc.}
\end{equation}

\medskip

\noindent\textbf{Definition (Definite BSE).}
The BSE eigenvalue problem~\eqref{eq:bse} is called \emph{definite}
if $\Hhat$ is positive definite.  This holds when the Tamm--Dancoff
approximation gives all-positive eigenvalues (physically, when the
system is stable against excitonic instabilities).  We assume the
definite case in what follows.


%======================================================================
\section{Reduction to a Hermitian eigenvalue problem}
\label{sec:hermitian}
%======================================================================

Since $\Hhat$ is Hermitian positive definite, it admits a Cholesky
factorization
\begin{equation}\label{eq:cholesky}
  \Hhat = L L^\hc
\end{equation}
where $L$ is lower triangular with positive diagonal entries.

\begin{theorem}[Hermitian reduction]\label{thm:hermitian}
  Define
  \begin{equation}\label{eq:K}
    K := L^\hc S L.
  \end{equation}
  Then $K$ is Hermitian and has the same eigenvalues as the BSE
  Hamiltonian $H$.
\end{theorem}

\begin{proof}
  \textbf{Hermiticity:}
  \[
    K^\hc = (L^\hc S L)^\hc = L^\hc S^\hc L = L^\hc S L = K,
  \]
  since $S$ is Hermitian.

  \textbf{Same eigenvalues as $H$:} We show that $K$ and $H$ are
  similar.  Since $L$ is invertible:
  \begin{align*}
    L^{-\hc} K L^\hc
    &= L^{-\hc} (L^\hc S L) L^\hc \\
    &= S L L^\hc \\
    &= S \Hhat \\
    &= H.
  \end{align*}
  Therefore $K = L^\hc H L^{-\hc}$, so $K$ and $H$ are related by a
  similarity transformation via $L^\hc$.  In particular, they have
  identical eigenvalues.
\end{proof}

\medskip

Since $K$ is Hermitian, its eigenvalues are real.  But $H$ has the
spectral symmetry that eigenvalues come in pairs $\pm\omega_k$
(see Section~\ref{sec:spectral}), so the eigenvalues of $K$ are
$\omega_1, -\omega_1, \omega_2, -\omega_2, \ldots, \omega_n, -\omega_n$
with $\omega_k > 0$.

\subsection{Eigenvector back-transformation}
If $K y = \omega y$, then the corresponding BSE eigenvector is
\begin{equation}\label{eq:backtransform}
  \boxed{x = L^{-\hc} y.}
\end{equation}
This follows because $H x = L^{-\hc} K L^\hc x = L^{-\hc} K y
= \omega L^{-\hc} y = \omega x$.  In practice, $x$ is obtained by
solving the triangular system $L^\hc x = y$.


%======================================================================
\section{Spectral symmetry and left eigenvectors}
\label{sec:spectral}
%======================================================================

\begin{theorem}[Spectral symmetry]\label{thm:spectral}
  If $\omega$ is an eigenvalue of $H$ with right eigenvector
  $x = (x_1, x_2)^T$, then $-\omega$ is also an eigenvalue with
  right eigenvector $(\conj{x}_2, \conj{x}_1)^T$.
\end{theorem}

\begin{proof}
  Let $Hx = \omega x$ with $x = (x_1, x_2)^T$.  Then:
  \begin{align*}
    Ax_1 + Bx_2 &= \omega x_1, \\
    -B^*x_1 - A^*x_2 &= \omega x_2.
  \end{align*}
  Take the complex conjugate of both equations:
  \begin{align*}
    A^* \conj{x}_1 + B^* \conj{x}_2 &= \conj{\omega}\, \conj{x}_1, \\
    -B \conj{x}_1 - A \conj{x}_2 &= \conj{\omega}\, \conj{x}_2.
  \end{align*}
  (We used $\conj{A} = A^T = A^{*T}$ and $\conj{B} = B^*$.)
  Interchanging the two equations and relabeling
  $\tilde{x}_1 = \conj{x}_2$ and $\tilde{x}_2 = \conj{x}_1$:
  \begin{align*}
    A \tilde{x}_1 + B \tilde{x}_2 &= -\conj{\omega}\, \tilde{x}_1, \\
    -B^* \tilde{x}_1 - A^* \tilde{x}_2 &= -\conj{\omega}\, \tilde{x}_2.
  \end{align*}
  This says $H \tilde{x} = -\conj{\omega}\, \tilde{x}$.  For the
  definite case, all eigenvalues are real ($\omega = \conj{\omega}$),
  so $-\conj{\omega} = -\omega$.
\end{proof}

\begin{theorem}[Left eigenvectors from right eigenvectors]%
  \label{thm:left}
  Let $Hx = \omega x$ with $x = (x_1, x_2)^T$.  Then
  $w = Sx = (x_1, -x_2)^T$ is a \textbf{left eigenvector} of $H$
  for the same eigenvalue $\omega$:
  \begin{equation}\label{eq:left}
    \boxed{w^\hc H = \omega w^\hc.}
  \end{equation}
\end{theorem}

\begin{proof}
  Since $H = S\Hhat$ and $S^2 = I$:
  \[
    w^\hc H = x^\hc S \cdot S \Hhat = x^\hc \Hhat.
  \]
  Also, from $Hx = \omega x$, i.e.\ $S\Hhat x = \omega x$, we get
  $\Hhat x = \omega S x$ (multiply by $S$).  Therefore:
  \[
    x^\hc \Hhat = (\Hhat x)^\hc = (\omega S x)^\hc = \omega x^\hc S = \omega w^\hc.
  \]
  (We used the fact that $\Hhat$ is Hermitian: $x^\hc \Hhat
  = (\Hhat x)^\hc$, and that $\omega$ is real in the definite case.)
\end{proof}

\begin{corollary}
  Left eigenvectors are obtained from right eigenvectors at no
  additional computational cost.  If the BSE states are ordered so
  that the first $n$ are resonant (transition energy $\Delta\epsilon > 0$)
  and the last $n$ are anti-resonant ($\Delta\epsilon < 0$), then the
  left eigenvector is obtained by negating the anti-resonant
  components.

  In GPAW's implementation, where states are ordered by
  $(k, v, c)$ indices, this generalizes to: multiply each component of
  the right eigenvector by $\operatorname{sign}(\Delta\epsilon_i)$.
\end{corollary}


%======================================================================
\section{Reduction to a real skew-symmetric eigenvalue problem}
\label{sec:skewsymm}
%======================================================================

When the BSE matrices $A$ and $B$ are \emph{real} (which occurs under
time-reversal symmetry for non-magnetic systems), a further reduction
is possible.

\begin{theorem}[Skew-symmetric reduction]\label{thm:skew}
  Let $H$ be a real BSE matrix ($A$, $B$ real symmetric).
  Let $\Hhat = LL^T$ be the (real) Cholesky factorization.  Define
  \begin{equation}
    K_J := L^T J L, \qquad J = \bmat{0 & I_n \\ -I_n & 0}.
  \end{equation}
  Then:
  \begin{enumerate}
    \item $K_J$ is real and skew-symmetric: $K_J^T = -K_J$.
    \item The eigenvalues of $K_J$ are $\pm i\omega_k$, where
      $\pm\omega_k$ are the eigenvalues of $H$.
  \end{enumerate}
\end{theorem}

\begin{proof}\leavevmode
  \begin{enumerate}
    \item Skew-symmetry:
      \[
        K_J^T = (L^T J L)^T = L^T J^T L = -L^T J L = -K_J.
      \]

    \item For the eigenvalue relation, we show that $K_J$ is similar
      to $J\Hhat$.  Since $L$ is invertible:
      \[
        L^{-T} K_J L^T = L^{-T}(L^T J L) L^T = J L L^T = J \Hhat.
      \]
      So $K_J$ and $J\Hhat$ have the same eigenvalues.

      Now we relate the eigenvalues of $J\Hhat$ to those of $H = S\Hhat$.
      Consider a simple case to establish the pattern: for $n = 1$
      with $A = a$ and $B = b$ (scalars),
      \[
        H = \bmat{a & b \\ -b & -a}, \quad
        J\Hhat = \bmat{0 & 1 \\ -1 & 0}\bmat{a & b \\ b & a}
        = \bmat{b & a \\ -a & -b}.
      \]
      The eigenvalues of $H$ are $\pm\sqrt{a^2 - b^2}$ and those of
      $J\Hhat$ are $\pm i\sqrt{a^2 - b^2}$.

      For the general case, note that both $H^2 = S\Hhat S\Hhat$ and
      $(J\Hhat)^2 = J\Hhat J\Hhat$ share a common structure.  Using
      $J\Hhat J = -S\Hhat S$ (which follows from direct computation),
      we obtain
      \[
        (J\Hhat)^2 = J \Hhat J \Hhat = -(S\Hhat S)\Hhat = -S\Hhat(S\Hhat) = -H^2.
      \]
      (The last step uses $S\Hhat = H$.)

      Therefore, if $\omega^2$ is an eigenvalue of $H^2$, then
      $-\omega^2$ is an eigenvalue of $(J\Hhat)^2$, which means
      $\pm i\omega$ are eigenvalues of $J\Hhat$ (and hence of $K_J$).
      \qedhere
  \end{enumerate}
\end{proof}

\medskip

\noindent\textbf{Verification of} $J\Hhat J = -S\Hhat S$:
\begin{align*}
  J\Hhat J
  &= \bmat{0 & I \\ -I & 0}\bmat{A & B \\ B & A}\bmat{0 & I \\ -I & 0}
  = \bmat{B & A \\ -A & -B}\bmat{0 & I \\ -I & 0} \\
  &= \bmat{-A & B \\ B & -A}.
\end{align*}
\begin{align*}
  S\Hhat S
  &= \bmat{I & 0 \\ 0 & -I}\bmat{A & B \\ B & A}\bmat{I & 0 \\ 0 & -I}
  = \bmat{A & -B \\ -B & A}.
\end{align*}
Indeed $J\Hhat J = -S\Hhat S$. \checkmark

\subsection{Advantages of the skew-symmetric formulation}

For a $2n \times 2n$ real skew-symmetric matrix:
\begin{itemize}
  \item Eigenvalues come in pairs $\pm i\sigma_k$, so only $n$ values
    need to be computed (factor-of-2 savings over a general solver).
  \item ELPA provides an efficient parallel solver
    (\texttt{elpa\_skew\_eigenvectors\_double}) that exploits a
    connection between skew-symmetric and symmetric tridiagonal forms,
    achieving up to $3.67\times$ speedup over reference ScaLAPACK
    implementations~\cite{Penke2020}.
  \item Memory usage is halved compared to a general complex solver.
\end{itemize}


%======================================================================
\section{Summary of the algorithm}
\label{sec:algorithm}
%======================================================================

\subsection{Complex BSE (general case)}
Given the BSE Hamiltonian $H = \bmat{A & B \\ -B^* & -A^*}$:

\begin{enumerate}
  \item Let $s_i = \operatorname{sign}(\Delta\epsilon_i)$ for each
    pair state $i$.  Form $S = \operatorname{diag}(s_1, \ldots, s_{2n})$.
  \item Compute the Hermitian companion $\Hhat = S H$.
  \item Cholesky factorize: $\Hhat = L L^\hc$.
  \item Form $K = L^\hc S L$ (Hermitian).
  \item Diagonalize: $K y_k = \omega_k y_k$ using a Hermitian
    eigensolver (ELPA or LAPACK).
  \item Back-transform right eigenvectors: $x_k = L^{-\hc} y_k$
    (solve $L^\hc x_k = y_k$).
  \item Left eigenvectors: $w_k = S x_k$ (negate anti-resonant
    components).
\end{enumerate}

Cost: $O(n^3)$ for Cholesky + $O(n^3)$ for eigensolver + $O(n^2)$ for
back-transformation.  The eigenvalues are guaranteed real and come in
$\pm\omega_k$ pairs.

\subsection{Real BSE (time-reversal symmetric case)}
When $A$ and $B$ are real symmetric:

\begin{enumerate}
  \item Compute $\Hhat = \bmat{A & B \\ B & A}$ (real symmetric
    positive definite).
  \item Real Cholesky: $\Hhat = L L^T$.
  \item Form $K_J = L^T J L$ (real skew-symmetric).
  \item Diagonalize with ELPA's skew-symmetric solver:
    eigenvalues $\sigma_k$ (the $\omega_k = \sigma_k$).
  \item Recover BSE eigenvectors from skew-symmetric eigenvectors.
\end{enumerate}

This gives a factor-of-2 savings over the complex Hermitian approach
because only $n$ eigenvalues (not $2n$) are computed.


\begin{thebibliography}{9}
\bibitem{Penke2020}
  C.~Penke, A.~Marek, C.~Vorwerk, C.~Draxl, and P.~Benner,
  ``High performance solution of skew-symmetric eigenvalue problems
  with applications in solving the Bethe--Salpeter eigenvalue problem,''
  \emph{Parallel Computing} \textbf{96}, 102639 (2020).
  \href{https://doi.org/10.1016/j.parco.2020.102639}%
  {DOI: 10.1016/j.parco.2020.102639}.

\bibitem{Shao2016}
  M.~Shao, F.~H.~da~Jornada, C.~Yang, J.~Deslippe, and S.~G.~Louie,
  ``Structure preserving parallel algorithms for solving the
  Bethe--Salpeter eigenvalue problem,''
  \emph{Linear Algebra Appl.} \textbf{488}, 148--167 (2016).
  \href{https://doi.org/10.1016/j.laa.2015.09.036}%
  {DOI: 10.1016/j.laa.2015.09.036}.

\bibitem{Benner2020}
  P.~Benner and C.~Penke,
  ``Efficient and accurate algorithms for solving the Bethe--Salpeter
  eigenvalue problem for crystalline systems,''
  \emph{J.\ Comput.\ Appl.\ Math.} \textbf{400}, 113650 (2022).
  \href{https://doi.org/10.1016/j.cam.2021.113650}%
  {DOI: 10.1016/j.cam.2021.113650}.
\end{thebibliography}

\end{document}
